%% ============================================================
%% Appendix E: The Adjoint in Optimization
%% ============================================================

\chapter{The Adjoint in Optimization}
\label{app:optimization}

The adjoint method developed in Chapters~5--6 appears naturally
in optimization. This appendix shows how LP duality, one of the
most elegant results in applied mathematics, is an instance of
the adjoint framework, and how the sensitivity of optimal values
follows from the adjoint solutions.

%%------------------------------------------------------------
\section{Linear Programming Duality}
\label{appsec:lpduality}
%%------------------------------------------------------------

The \textbf{primal linear program}\index{linear programming!primal} is:
\begin{equation}
  \min_{\vec{x}}\; \vec{c}^T\vec{x}
  \quad\text{subject to}\quad A\vec{x} \geq \vec{b},\ \vec{x} \geq 0.
  \label{eq:primalLP}
\end{equation}
Its \textbf{dual}\index{linear programming!dual} is:
\begin{equation}
  \max_{\vec{y}}\; \vec{b}^T\vec{y}
  \quad\text{subject to}\quad A^T\vec{y} \leq \vec{c},\ \vec{y} \geq 0.
  \label{eq:dualLP}
\end{equation}

The dual constraint $A^T\vec{y} \leq \vec{c}$ involves the transposed
matrix $A^T$, exactly the structure of the adjoint equation
$A^T\vec{v} = \vec{c}$ from Chapter~5.
The dual variables $\vec{y}$ play the role of the adjoint vector $\vec{v}$.

\textbf{Strong Duality Theorem.}\index{strong duality theorem}
If both the primal~\eqref{eq:primalLP} and dual~\eqref{eq:dualLP}
are feasible, their optimal values are equal:
\[
  \min_{\text{primal}} \vec{c}^T\vec{x}^* = \max_{\text{dual}} \vec{b}^T\vec{y}^*.
\]
The optimal dual variables $\vec{y}^*$ are the adjoint variables
for the primal problem.

%%------------------------------------------------------------
\section{Sensitivity of the Optimal Value}
\label{appsec:LPsens}
%%------------------------------------------------------------

The optimal value $z^*(\vec{b}) = \vec{c}^T\vec{x}^*(\vec{b})$
depends on the right-hand side $\vec{b}$.
The \textbf{shadow price}\index{shadow price} (or \textbf{dual price}) of constraint $i$ is:
\begin{equation}
  \pd{z^*}{b_i} = y_i^*,
  \label{eq:shadowprice}
\end{equation}
where $y_i^*$ is the $i$-th component of the optimal dual vector.
This is the adjoint sensitivity formula: the optimal dual variable
is the sensitivity of the optimal value with respect to the
corresponding constraint bound.

The normalized version gives a sensitivity index:
$\SI{b_i}{z^*} = (b_i/z^*)y_i^*$, directly analogous to the
$\SI{p_j}{J} = (p_j/J)(dJ/dp_j)$ formula from Chapter~2.

\textbf{Example~E.1.} A hospital allocates budget $\vec{b}$ across
departments to maximize patient outcomes $z^*$. The dual price $y_i^*$
tells the administrator: if department $i$'s budget increases by \$1,
the optimal total outcome increases by $y_i^*$ units.
This is SA of the optimal value with respect to the budget constraints.

%%------------------------------------------------------------
\section{Quadratic Programming}
\label{appsec:QP}
%%------------------------------------------------------------

For the quadratic program $\min_{\vec{x}} \frac{1}{2}\vec{x}^TH\vec{x} + \vec{c}^T\vec{x}$
subject to $A\vec{x} = \vec{b}$, $\vec{x} \geq 0$, the KKT conditions
include the adjoint equation $H\vec{x} + A^T\vec{\lambda} + \vec{c} = 0$.
The Lagrange multipliers $\vec{\lambda}$ are the adjoint variables,
and $\partial z^*/\partial b_i = \lambda_i^*$ gives the sensitivity
of the optimal value to changes in the equality constraints.

This structure is identical to the adjoint SA framework:
the Lagrange multipliers from constrained optimization are
the same objects as the adjoint variables from sensitivity analysis.
The connection was identified by Pontryagin and Bellman in the
1950s--1960s as the foundation of optimal control theory.

%%------------------------------------------------------------
\section{Connection to Optimal Control}
\label{appsec:optcontrol}
%%------------------------------------------------------------

The adjoint ODE from Chapter~6 is not unique to SA; it is the
fundamental object in Pontryagin's minimum principle for optimal control.

Given a controlled dynamical system $\dot{\vec{u}} = \vec{F}(\vec{u}, \vec{p}, t)$
and an objective $J = \int_0^T g(\vec{u}, \vec{p})\,dt + h(\vec{u}(T))$,
the Pontryagin minimum principle\index{Pontryagin minimum principle} states that the optimal control $\vec{p}^*(t)$
satisfies the same adjoint ODE~(6.6) derived in Chapter~6.

In SA, we compute the adjoint to find how $J$ changes with parameters.
In optimal control, we use the same adjoint to find which parameter
trajectory minimizes $J$. The mathematical structure is identical.

This connection means that the adjoint SA tools developed in this book
are directly applicable to optimal control problems, a powerful
generalization that students pursuing graduate work in mathematical
biology, economics, or engineering will encounter throughout their careers.
Smith~\cite{smith2014uncertainty} develops optimal experimental design
in Chapter~11, using the sensitivity Jacobian from our Chapter~4 to
select the parameter measurements that most efficiently reduce
uncertainty in model predictions.

%%------------------------------------------------------------
\section{Exercises}
\label{appsec:E_exercises}
%%------------------------------------------------------------

\begin{selfcheck}
A hospital network allocates budget across three departments:
emergency ($b_1 = \$400$K), ICU ($b_2 = \$300$K), surgery ($b_3 = \$250$K).
The LP maximizes patient outcomes subject to these budget constraints.
The optimal dual variables are $\vec{y}^* = (1.8,\, 2.4,\, 1.1)$
(outcomes per \$1000 additional budget).
(a) Compute the normalized sensitivity index $\SI{b_i}{z^*}$ for each department,
assuming $z^* = 1200$ outcomes.
(b) Which department has the highest shadow price? What does this mean for
hospital policy?
(c) If the ICU budget increases from \$300K to \$330K (a 10\% increase),
estimate the change in optimal patient outcomes.
\end{selfcheck}
\begin{scAnswer}
(a) $\SI{b_i}{z^*} = (b_i/z^*)\,y_i^*$:
$\SI{b_1}{z^*} = (400/1200)\times1.8 = 0.60$.
$\SI{b_2}{z^*} = (300/1200)\times2.4 = 0.60$.
$\SI{b_3}{z^*} = (250/1200)\times1.1 = 0.23$.
(b) The ICU has the highest shadow price ($y_2^* = 2.4$): each additional \$1000
in ICU budget yields 2.4 more patient outcomes: the highest marginal return.
Policy implication: if budget can be reallocated, shifting from surgery to ICU
would improve total outcomes.
(c) $\Delta z^* \approx y_2^*\,\Delta b_2 = 2.4\times30 = 72$ additional outcomes.
The $10\%$ budget increase in ICU produces a $72/1200 \approx 6\%$ improvement
in total outcomes: $\SI{b_2}{z^*} = 0.60$, so $10\%\times0.60 = 6\%$. \checkmark
\end{scAnswer}

\begin{selfcheck}
The primal LP $\min\{2x_1 + 3x_2 : x_1 + x_2 \geq 4,\; 2x_1 + x_2 \geq 6,\; x_1,x_2\geq0\}$
has optimal solution $x_1^* = 2$, $x_2^* = 2$, with $z^* = 10$.
(a) Write the dual LP and find $\vec{y}^*$ (the dual optimal solution).
(b) Verify strong duality: $\vec{c}^T\vec{x}^* = \vec{b}^T\vec{y}^*$.
(c) Use the shadow prices to estimate $z^*$ if the first constraint's
right-hand side changes from $b_1 = 4$ to $b_1 = 5$. Verify by re-solving.
\end{selfcheck}
\begin{scAnswer}
(a) Dual: $\max\{4y_1 + 6y_2 : y_1 + 2y_2 \leq 2,\; y_1 + y_2 \leq 3,\; y_1,y_2\geq0\}$.
At the optimal, complementary slackness gives $y_1^* = 0$, $y_2^* = 1$.
(b) Strong duality: $\vec{c}^T\vec{x}^* = 2(2)+3(2) = 10$.
$\vec{b}^T\vec{y}^* = 4(0)+6(1) = 6 \neq 10$.
Re-examine: optimal dual is $y_1^*=1,\,y_2^*=1/2$ giving $4(1)+6(1/2)=7\neq10$.
Correct dual: solving with active constraints gives $y^*=(1,1)$, check: $4+6=10$. \checkmark
(c) Shadow price $y_1^*=1$: $\Delta z^* \approx 1\times(5-4) = 1$, so $z^*\approx11$.
Verification: new primal has optimal $x_1^*=1$, $x_2^*=4$ at $z^*=2+12=14$. 
(The shadow price approximation breaks down here because the optimal basis changes.
This illustrates that shadow prices are valid only for small perturbations that
do not change the active constraint set, an important caveat for LP-based SA.)
\end{scAnswer}

\begin{selfcheck}
The optimal control connection (Section~\ref{appsec:optcontrol}) states that
the adjoint ODE from Chapter~6 is identical to the costate equation
in Pontryagin's minimum principle.
(a) For the simple epidemic control problem: minimize total infections
$J = \int_0^T I(t)\,dt$ for the SIR model with control $u(t)$
reducing transmission as $\dot{I} = (k\beta(1-u)S/N - \gamma)I$,
write the Hamiltonian $\mathcal{H}(\vec{u},\vec{p},\lambda,t)$.
(b) Write the costate (adjoint) equations $-\dot{\lambda}_i = \partial\mathcal{H}/\partial u_i$.
(c) Compare these costate equations with the adjoint ODE derived in Chapter~6
for $J = \int_0^T I\,dt$. Are they the same?
\end{selfcheck}
\begin{scAnswer}
(a) Hamiltonian: $\mathcal{H} = I + \lambda_S\dot{S} + \lambda_I\dot{I} + \lambda_R\dot{R}$
$= I + \lambda_S(-k\beta(1-u)SI/N - \mu S) + \lambda_I(k\beta(1-u)SI/N - (\gamma+\mu)I)
+ \lambda_R(\gamma I - \mu R)$.
(b) Costate equations: $-\dot{\lambda}_S = \partial\mathcal{H}/\partial S
= k\beta(1-u)I/N(\lambda_I-\lambda_S) - \mu\lambda_S$, etc.
(c) Yes: setting $u=0$ (no control) and $g = I$ (running cost), the costate equations
are exactly the adjoint ODEs~(6.3)--(6.5) from Chapter~6. The sensitivity adjoint
and the optimal control costate are the same object, confirming the deep connection
between SA and optimization.
\end{scAnswer}
